\subsection{Legal and ethical implications of AI systems gaining personhood}
\label{sec:8.4.4}
Whether an AI system should ever hold personhood is worked out in sections 8.6.1 and 8.6.2. What belongs in a chapter about policy frameworks is a narrower question those sections do not settle: who decides, and when.

Every version of this question that has actually reached a legislature so far arrived as a response to pressure — a liability dispute, a public controversy, a product already on the market — rather than from a standing body asking it in advance. That ordering is the problem. A legislature deciding under pressure decides badly, and it decides against a factual record assembled by whoever had an interest in assembling one. The interdisciplinary review body section 8.6.1 describes only does the work it is designed for if some jurisdiction creates it before a forcing case shows up, which means creating it while the question still looks academic and nobody's damages are riding on the answer.

That is a policy choice rather than a doctrinal one, and it is cheap relative to almost everything else in this chapter. A national AI strategy that names an agency, gives it a mandate to keep the question under review, and sets a trigger for revisiting it — a capability threshold, a deployment scale, a category of use — decides in advance who will be in the room. It does not decide the answer, and it should not try to. It decides that when the answer is needed there will be a body qualified to give one, which is the part no forcing case ever supplies on its own.
